\section{Connections to Related Work}\label{sec:related}

Our focus on local merge compatibility draws on, and extends, several established theoretical traditions. We view our framework as a semantic generalization of randomized sketching, justified by decision-theoretic sufficiency.

\paragraph{Mergeable Summaries and Sketches.} 
The closest formal analogue to our work is the theory of \emph{mergeable summaries} for data streams \citep{Agarwal2013MergeableTODS}. In that setting, a summary (or sketch) $S$ supports a binary merge operator $\oplus$ such that $S(A \cup B) \approx S(A) \oplus S(B)$, with error guarantees that hold under arbitrary merge trees. 
Our framework can be understood as generalizing this logic to the semantic, non-commutative domain of natural language.
Specifically, when the task oracle $\fstar$ is symmetric (e.g., word counts or bag-of-words statistics) and the metric is Euclidean, our \emph{Merge Consistency} condition (C2) is mathematically isomorphic to the definition of a mergeable summary. In this specific case, checking Condition~\eqref{eq:leaf_compatibility} is equivalent to verifying that the LLM acts as a valid adder for a Count-Min sketch \citep{CormodeMuthukrishnan2005}.
However, text data is rarely commutative---the order of paragraphs changes the meaning of a narrative. Our contribution is to impose the rigorous discipline of mergeable summaries on these non-commutative operations, replacing the explicit algebra of sketches with ``learned'' local consistency checks that force a stochastic summarizer to behave \emph{as if} it were a mergeable data structure.
Appendix~\ref{app:mergeable} formalizes this bridge, mapping our notation $(\Strings,\concat,\fstar,g)$ onto the streaming literature and proving that Condition~\eqref{eq:leaf_compatibility} coincides with mergeability when the oracle is symmetric.

\paragraph{Decision Theory and Blackwell Sufficiency.}
A second foundation is decision theory. Under proper scoring rules, the optimal Bayesian action depends only on the posterior of the target; reporting the truth is optimal \citep{GneitingRaftery2007,Dawid2007}. Our Assumption~\ref{ass:CF} instantiates this principle: if supervision factorizes through $\fstar(X)$, then $\fstar$ is a \emph{sufficient statistic} for the population decision problem. Replacing $X$ with a summary that preserves $\fstar$ leaves the Bayes risk unchanged. This is a task-specific version of Blackwell sufficiency \citep{Blackwell1953,Torgersen1991}, where $\fstar$ acts as the sufficient reduction. While Blackwell's theorem is often invoked abstractly, our framework provides an operational audit to verify that a specific summarization pipeline actually retains this sufficiency in practice.

\paragraph{Information-theoretic parallels.}
Our framework is grounded in sufficiency and mergeable summaries, but it also admits an information-theoretic interpretation: the summary acts as the sufficient component while the discarded text serves as an ancillary ``remainder.'' For readers interested in that lens, Appendix~\ref{sec:info-sieve} sketches how the decomposition ideas behind the \emph{Information Sieve} of \citet{steeg2016informationsieve} can be adapted to oracle-preserving compression, though this viewpoint is entirely optional for the main results.

\paragraph{Preference Learning and RLHF.}
Modern alignment techniques, such as Direct Preference Optimization (DPO), model preferences using Bradley–Terry–Luce likelihoods driven by a reward function \citep{BradleyTerry1952,RafailovEtAl2023DPO}. A standard assumption in this literature is that human preferences depend on the input only through its task-relevant semantics (the oracle value), not incidental phrasing. Under this factorization, our preservation results imply that training DPO on oracle-preserving summaries attains the same population optimum as training on full documents. Absent such preservation, empirical performance gains can be illusory: a learned policy might entangle with artifacts of a particular merge schedule---a failure mode our audit is specifically designed to detect. When readability or fluency preferences fall outside $\fstar$, they must be modeled by a richer oracle or handled by separate alignment channels.

\paragraph{Global Structure vs. Local Checks.}
Finally, we contrast our approach with corpus-level abstractions like Latent Dirichlet Allocation (LDA) or Structural Topic Models (STM) \citep{BleiNgJordan2003,RobertsStewartTingley2014}. While those methods discover latent global structures unsupervised, we assume a pre-defined outcome label $\fstar$ determined by the researcher.
Furthermore, while our results induce a monoid congruence on strings when assumptions hold globally (discussed in Section~\ref{sec:global}), we deliberately avoid assuming a global algebra. Instead, we rely on ``local'' consistency: checks enforced only on the specific leaves and merges realized by the scheduler. This allows us to estimate the fidelity of a pipeline without requiring the oracle space $\Y$ to possess a strict algebraic structure that the task semantics may not justify. The tradeoff is scope: we restrict attention to decomposable coding, extraction, counting, or guided-summarization tasks where $\fstar$ is defined on bounded spans.

\bigskip
\noindent
In short, we borrow the algebraic intuition of mergeable sketches, adopt the sufficiency criteria of decision theory, and insist on edge-wise, testable equalities where practitioners actually operate. With this context in place, we turn to the formal setup.
